\documentclass[10pt]{article}
\usepackage[utf8]{inputenc}
\usepackage[T1]{fontenc}
\usepackage{amsmath}
\usepackage{amsfonts}
\usepackage{amssymb}
\usepackage[version=4]{mhchem}
\usepackage{stmaryrd}
\usepackage{graphicx}
\usepackage[export]{adjustbox}
\graphicspath{ {./images/} }
\usepackage{hyperref}
\hypersetup{colorlinks=true, linkcolor=blue, filecolor=magenta, urlcolor=cyan,}
\urlstyle{same}

\DeclareUnicodeCharacter{2297}{\ifmmode\otimes\else{$\otimes$}\fi}
\DeclareUnicodeCharacter{2713}{\ifmmode\checkmark\else{$\checkmark$}\fi}
\DeclareUnicodeCharacter{25EF}{\ifmmode\bigcirc\else{$\bigcirc$}\fi}

\begin{document}
\section*{Factor a polynomial Root into Roots of smallest possible degree}
Asked 10 years, 7 months ago Modified 7 years ago Viewed 700 times

\begin{itemize}
  \item[] Suppose I have a polynomial Root representing an algebraic number. I want to represent it as a product of several polynomial Root s (if possible) such that the largest degree among the factors is as small as possible. For example, $\operatorname{Root}\left[-1-\# 1+3 \# 1^{\wedge} 3-\# 1^{\wedge} 4+\# 1^{\wedge} 5-3 \# 1^{\wedge} 6+2 \# 1^{\wedge} 7+\right.$ \#1\^{}9 \&, 1] has to be represented as $\operatorname{Root}\left[1+\#+\#^{\wedge} 3\right.$ \&, 1] Root[1 - \# + \#\^{}3 \&, 1].
  \item[] The second problem is the same, but replacing "product" with "sum". For example, Root[8-4 \#1 $\left.+24 \# 1^{\wedge} 2-15 \# 1^{\wedge} 3+3 \# 1^{\wedge} 5+6 \# 1^{\wedge} 6+\# 1^{\wedge} 9 \&, 1\right]$ has to be represented as $\operatorname{Root}\left[1+\#+\#^{\wedge} 3\right.$ \&, 1] + Root[1 - \# + \#\^{}3 \&, 1].
\end{itemize}

How can I implement solutions to these problems in Mathematica?

\begin{center}
\begin{tabular}[t]{|l|l|l|l|l|l|}
\hline
simplifying-expressions & symbolic & algebraic-manipulation & algorithm & number-theory &  \\
\hline
Share Edit Close Delete Flag &  & edited Aug 7, 2019 at 18:45 &  & \begin{tabular}[t]{l}
asked Feb 8, 2016 at 18:55\includegraphics[max width=\textwidth, alt={}]{0211096f-960b-436d-8f9d-0de2165ff460-1_88_88_1542_1450}
 \\
6,651 \\
\end{tabular} & 3283 \\
\hline
\end{tabular}
\end{center}

\begin{itemize}
  \item[5] I'm sure you are aware of this but I'll say it anyway: this is not an easy problem. - Daniel Lichtblau May
  \begin{itemize}
    \item[] 9, 2016 at 14:14
    \item[-] Would this perhaps be better asked on Mathematics? - Mr.Wizard Jul 8, 2016 at 6:31
    \item[F] 
    \item[-] Recently we worked on a Metropolis-Hasting algorithm here in permutation-98 space. Consider two polynomials in rational coefficients of degree 3. That's rational-8 space. Is it inconceivable to adapt this algorithm to search this space and converge to a solution? - Dominic Aug 7, 2019 at 20:49
    \item[1] @Dominic for Metropolis-Hastings you need to define a reasonably continuous quality function that you can try to maximize or minimize. Do you have any ideas for how to define it? - Roman Aug 7, 2019 at 22:57
  \end{itemize}
  \item[2] Just to state the obvious, going in opposite direction can be done with the help of direct products of - companion matrices. Here, the problem is to write a companion matrix as a direct product of companion matrices. \href{http://math.stackexchange.com/questions/331017/}{math.stackexchange.com/questions/331017/}... - yarchik Aug 8, 2019 at 7:35
\end{itemize}

\section*{A constructive approach}
13 The problem can be solved if the form of the solution is given.

\section*{Define the two factors using a hint (that these should be cubic equations) in the original post}
$\mathrm{y} 1=\mathrm{x}^{\wedge} 3-\mathrm{p} 1 \mathrm{x}+\mathrm{q} 1$

\begin{verbatim}
y2 = x^3 - p2 x + q2
\end{verbatim}

+500 Build a companion matrix of the polynomial $p(x)$

\begin{verbatim}
CompanionMatrix[p_,x_]:=Module[{n,w=CoefficientList[p,x]},w=-w/Last[w];
    n=Length[w]-1; SparseArray[{{i_,n}:>w[[i]],{i_,j_}/;i==j+1->1},{n,n}]]
\end{verbatim}

\begin{itemize}
  \item[-] The roots of a polynomial equation $p_{A}(x)=0$ are given by the eigenvalues of its companion matrix $A$.
  \item[-] If $a$ is a root of $p_{A}(x)$ (with companion matrix $A$ ) and $b$ is a root of $p_{B}(x)$ (with companion matrix $B$ ) then $a b$ is an eigenvalue of $A \otimes B$ and $a+b$ is an eigenvalue of $A \otimes I+I \otimes B$, where ⊗ is the direct product of matrices.
\end{itemize}

Let us focus on the product case and, therefore, determine the characteristic equation of the direct product of companion matrices

\begin{verbatim}
(a1=CompanionMatrix[y1,x])//MatrixForm
(a2=CompanionMatrix[y2,x])//MatrixForm
y3=CharacteristicPolynomial[KroneckerProduct[a1,a2],x]
y4=y3 Sign[CoefficientList[y3,x]//Last]
\end{verbatim}

The sum is treated similarly.\\
We find out that the resulting characteristic polynomial is

\begin{verbatim}
y4=-q3^3 q2^3 + p1 p2 q1^2 q2^2 x + (-p2^3 q1^2 - p1^3 q2^2 + 3 q1^2 q2^2) x^3 +
p1 p2 q1 q2 x^4 + p1^2 p2^2 x^5 - 3 q1 q2 x^6 - 2 p1 p2 x^7 + x^9
\end{verbatim}

Notice the relatively simple form in the case of trinomial equations. Now, your polynomial is

\begin{verbatim}
z = -1 - #1 + 3 #1^3 - #1^4 + #1^5 - 3 #1^6 + 2 #1^7 + #1^9 &@x
\end{verbatim}

We demand that the two polynomials ( y3 and z ) are equal for any value of x

\begin{verbatim}
r = SolveAlways[z == y4, {x}]
\end{verbatim}

Several solutions are obtained. Take, for instance, the first one

\begin{verbatim}
{y1, y2} /. r[[1]]
(*{1/q2 + x/q2^(2/3) + x^3, q2 - q2^(2/3) x + x^3}*)
\end{verbatim}

Setting $\mathrm{q} 2=1$ we obtain the OP result.

\section*{Comment on the method}
MA has a nice RootApproximant function, which operates by virtue of the LLL algorithm. One may be tempted to follow this route and implement some kind of this experimental mathematics approach. In contrast, the presented solution is fully constructive (in fact, algebraic) and does not require arbitrary precision computations.

\section*{Answer to the 1st challenge question}
\begin{verbatim}
y1=x^5+ m x^3+n x^2+p1 x+q1;
y2=x^3+ p2 x+q2;
a1=CompanionMatrix[y1,x];
a2=CompanionMatrix[y2,x];
y3=CharacteristicPolynomial[KroneckerProduct[a1,a2],x];
y4=y3 Sign[CoefficientList[y3,x]//Last];
z=-282300416-64550400 #-34426880 #^2-14185880 #^3+8564800 #^4+4231216 #^5-972800
#^6-367820 #^7+27360 #^8+2600 #^9+1680 #^10+100 #^11-240 #^12+40 #^13+#^15&@x;
r=SolveAlways[y4==z,{x}]
({y1,y2}/.r[[1]])/.q2->1
Out[1]= {656-150 x+80 x^2-20 x^3+x^5,1+x+x^3}
\end{verbatim}

\begin{itemize}
  \item[-] Notice that in the characteristic equation the highest order term can be negative, whereas I assume that in the given equation z it is always 1. Therefore, in order to use SolveAlways y3 is multiplied by the sign.
  \item[-] From the shown solutions, one can see that there is some arbitrariness in the results. But this is, of course, expected. As for the shown solution in radicals, one probably needs to guess the field extension.
  \item[-] One can create a Module as to fully automate the derivation. But is there a pressing need?
\end{itemize}

\section*{Answer to the 2st challenge question}
Here I demonstrate that the method can be used to write a root $P$ in terms of three roots of lower-order polynomials as $P=Q+R S$. Additionally, the computation is done in a more structured form

\begin{enumerate}
  \item Define a generic polynomial with 1 as the leading coefficient
\end{enumerate}

\begin{verbatim}
pX[a_,n_]:=x^n+Sum[a[i]x^i,{i,0,n-1}]
\end{verbatim}

\begin{enumerate}
  \setcounter{enumi}{1}
  \item Define two modules that split a root in terms of a sum or a product
\end{enumerate}

\begin{verbatim}
pSum[pA_,pB_,y_]:=Module[{mA,mB,mAB,nA,nB,p},
    nA=Exponent[pA,y];
    nB=Exponent[pB,y];
    mA=CompanionMatrix[pA,y];
    mB=CompanionMatrix[pB,y];
mAB=KroneckerProduct[mA,IdentityMatrix[nB]]+KroneckerProduct[IdentityMatrix[nA],mB];
    p=CharacteristicPolynomial[mAB,y];
    p Sign[CoefficientList[p,y]//Last]
]
pProduct[pA_,pB_,y_]:=Module[{mA,mB,mAB,p},
    mA=CompanionMatrix[pA,y];
    mB=CompanionMatrix[pB,y];
    mAB=KroneckerProduct[mA,mB];
    p=CharacteristicPolynomial[mAB,y];
    p Sign[CoefficientList[p,y]//Last]
]
\end{verbatim}

\begin{enumerate}
  \setcounter{enumi}{2}
  \item Construct a working example for polynomials of the degrees 3, 2, 2 respectively.
\end{enumerate}

\begin{verbatim}
{nQ,nR,nS}={3,2,2};
nT=nR nS;
pP[0]=(RootReduce[Root[#^3+#+3&,1]+Root[#^2+#+5&,1]Root[#^2+#+7&,1]]//First)@x
Out[1]= 1984051825-172403780 x-281288553 x^2+14148329 x^3+17544721 x^4-310509
x^5-619703 x^6-4623 x^7+13443 x^8+244 x^9-167 x^10-3 x^11+x^12
\end{verbatim}

\begin{enumerate}
  \setcounter{enumi}{3}
  \item Do the first part, namely, split the root of a 12th order equation into a sum of 3rd and 4th order roots
\end{enumerate}

\begin{verbatim}
(sol[1]=SolveAlways[pP[0]==pSum[pX[q,nQ],pX[t,nT],x],x])//Transpose//TableForm
rule[1]=q[2]->0;
sol[1,1]=First[sol[1]]/.rule[1];
AppendTo[sol[1,1],rule[1]];
{pX[q,nQ],pX[t,nT]}/.sol[1,1]
\end{verbatim}

\begin{verbatim}
Out[2]= {3+x+x^3,1225-35 x-58 x^2-x^3+x^4}
\end{verbatim}

\begin{itemize}
  \item[5.] Do the second part, split the 4th order root into a product of 2 roots of quadratic equations
\begin{verbatim}
(sol[2]=SolveAlways[CpX[t,nT]/.sol[1,1])==pProduct[pX[r,nR],pX[s,nS],x],x])//Transpc
rule[2]=s[1]->1;
sol[2,1]=First[sol[2]]/.rule[2];
AppendTo[sol[2,1],rule[2]];
{pX[r,nR],pX[s,nS]}/.sol[2,1]
Out[3]= {5+x+x^2,7+x+x^2}
\end{verbatim}


  \item[✓] 
  \item[6.] The coefficients in rule[1] and rule[2] have been selected as to match the original equation. However, other choices are possible.
\end{itemize}

\begin{center}
\includegraphics[max width=\textwidth, alt={}]{0211096f-960b-436d-8f9d-0de2165ff460-5_149_1600_1045_271}
\end{center}

1 Very nice (+ upvote, of course). I missed a beat on this one-- it can similarly be done using resultants

\begin{itemize}
  \item[] and solving. Basically equivalent to your approach, if you drill deep enough. - Daniel Lichtblau Aug 8, 2019 at 14:54
  \item[1] E.g. Root[-282300416 - 64550400 \# - 34426880 \#\^{}2 - 14185880 \#\^{}3 + 8564800 \#\^{}4 + 4231216 \#\^{}5 - 972800 \#\^{}6 - 367820 \#\^{}7 + 27360 \#\^{}8 + 2600 \#\^{}9 + 1680 \#\^{}10 + 100 \#\^{}11 - 240 \#\^{}12 + 40 \#\^{}13 + \#\^{}15 \&, 1] can be represented as Root[-656 - 150 \# - 80 \#\^{}2 - 20 \#\^{}3 + \#\^{}5 \&, 1] * Root[-1 + \# + \#\^{}3 \&, 1] and ... - Vladimir Reshetnikov Aug 8, 2019 at 18:38 /
  \item[1.] .. actually, can be expressed in radicals:
$$
\frac{\left(\sqrt[3]{6} \beta^{2}-23^{2 / 3}\right)\left(\alpha\left(\sqrt[10]{2} \gamma^{4}+\left(\sqrt[5]{2}+2^{7 / 10}\right) \gamma^{2}+2 \gamma+2^{3 / 10}(5+3 \sqrt{2})\right)+(-68-35 \sqrt{2}+2 \sqrt{4100-1918 \sqrt{2}}) \gamma-22^{3 / 10}(68+35 \sqrt{2})\right)}{32^{17 / 30} \alpha \beta \gamma^{3}},
$$
where $\alpha=\sqrt{2050-959 \sqrt{2}}, \beta=\sqrt[3]{9+\sqrt{93}}, \gamma=\sqrt[5]{102+72 \sqrt{2}-\sqrt{20714+14647 \sqrt{2}}}$
  \begin{itemize}
    \item[-] Vladimir Reshetnikov Aug 8, 2019 at 18:39
  \end{itemize}
  \item[1] @VladimirReshetnikov Try $\sqrt{21}$ as an extension. Yes, these were my guesses, more systematic way is - through the computation of the Galois group of the polynomial. There are some MA packages available (\href{http://library.wolfram.com/infocenter/TechNotes/158}{library.wolfram.com/infocenter/TechNotes/158}) and some nice posts on math.stackexchange: \href{http://math.stackexchange.com/questions/45893/}{math.stackexchange.com/questions/45893/}... . - yarchik Aug 11, 2019 at 19:54
  \item[1] Thanks! By the way, I tried the package Solving the Quintic with Mathematica you linked, and it worked unreliable for me, failing to found existing solutions sometimes. Here is my implementation of a method from Daniel Lazard's paper Solving Quintics by Radicals:\\
\href{http://mathematica.stackexchange.com/a/134551/7288}{mathematica.stackexchange.com/a/134551/7288}. I corrected a few typos found in the paper to make it work, and so far have not found any bugs in it. - Vladimir Reshetnikov Aug 11, 2019 at 20:24
\end{itemize}

I'll show the resultant formulation for the degree 15 example.

The polynomial in question:\\
poly15 $=\left(-282300416-64550400 \# \frac{14426880}{} \#^{\wedge} 2-14185880 \#^{\wedge} 3+\right.$

\begin{verbatim}
8564800 #^4 + 4231216 #^5 - 972800 #^6 - 367820 #^7 +
27360 #^8 + 2600 #^9 + 1680 #^10 + 100 #^11 - 240 #^12 +
40 #^13 + #^15) &[z];
\end{verbatim}

Define monic polynomials of degree 3 and 5 with symbolic coefficients.

\begin{verbatim}
poly3[x_] := Array[a, 3, 0].x^Range[0, 2] + x^3
poly5[x_] := Array[b, 5, 0].x^Range[0, 4] + x^5
\end{verbatim}

Here is the double-resultant way to get a polynomial whose roots are products of the roots of poly3 and poly5.

\begin{verbatim}
resxy = Resultant[Resultant[z - x*y, poly3[y], y], poly5[x], x];
\end{verbatim}

Equate coefficients with this and the degree 15 polynomial. Since we can scale the two roots, force a constant term to be unity.

\begin{verbatim}
coeffs = CoefficientList[resxy - poly15, z] /. {a[0] -> 1}
\end{verbatim}

Now solve:

\begin{verbatim}
solns=Solve[coeffs == 0]
(* Out[285]= {{a[1] -> 1, a[2] -> 0, b[0] -> 656, b[1] -> -150,
    b[2] -> 80, b[3] -> -20, b[4] -> 0}, {a[1] -> -(-1)^(1/3),
    a[2] -> 0, b[0] -> 656 (-1 + (-1)^(1/3)), b[1] -> 150 (-1)^(1/3),
    b[2] -> 80, b[3] -> 20 (1 - (-1)^(1/3)),
    b[4] -> 0}, {a[1] -> (-1)^(2/3), a[2] -> 0,
    b[0] -> 656 (-1 - (-1)^(2/3)), b[1] -> -150 (-1)^(2/3), b[2] -> 80,
    b[3] -> 20 (1 + (-1)^(2/3)), b[4] -> 0}} *)
\end{verbatim}

Recover polynomials that give the root pairs. To get integer coefficients we use the first solution above.

\begin{verbatim}
{poly3[x], poly5[x]} /. solns[[1]]
(* Out[289]= {x + x^3 + a[0], 656-150 x + 80 x^2 - 20 x^3 + x^5} *)
\end{verbatim}

Not so well explained as the approach by @yarchik (which I'd expect to see get the bounty), but at some level it is equivalent.

Share Edit Follow Flag\\
answered Aug 8, 2019 at 21:43

\begin{itemize}
  \item[◯] Nice trick with 2 Resultant s (+). I thought, you would be using 1 resultant. - yarchik Aug 9, 2019 at
  \item[◯] 11:12
  \item[◯] @yarchik It can be done with one but then one of the polynomials needs to have a variable
  \item[◯] substitution, something like $x-->y / x$, and I can never remember exactly what it is. - Daniel Lichtblau Aug 9, 2019 at 15:00
\end{itemize}

\end{document}